\section{Anamnesis, Offering, and Sacrifice}

The Greek noun \emph{anamnesis} occurs in the institution command in Luke
22:19 and 1 Corinthians 11:24--25 and in Hebrews' contrast between repeated
sacrifices and the annual remembrance of sins (Heb 10:3). Lexical range alone
cannot prove a complete Catholic theology. Context controls. At the Supper the
command is attached to Christ's Body given and covenant Blood poured out. In
Paul the action proclaims the Lord's death until he comes. The Church's
Eucharistic anamnesis is therefore enacted proclamation under Christ's command,
not merely an internal recollection.

The New Testament also controls how ``sacrifice'' is used. Christ hands himself
over in obedience and love (Rom 5:6--11; Gal 2:20; Eph 5:2, 25). Hebrews
presents him as high priest who offers himself and enters the heavenly
sanctuary, while excluding
repeated immolations. Revelation's Lamb is slain yet standing. The Church's
offering is derivative and participatory: believers offer their bodies as a
living sacrifice, praise, alms, and mutual service through Christ (Rom 12:1;
Phil 4:18; Heb 13:15--16; 1 Pet 2:5).

Three errors follow from confusing these registers. First, one can make
``memorial'' so mental that Christ's objective sacramental action disappears.
Second, one can make ``sacrifice'' mean another killing and thereby contradict
Hebrews. Third, one can make ecclesial offering compete with Christ's. The
Catholic synthesis avoids all three. Anamnesis is the Church's effective
liturgical memorial under Christ's command. Christ offered himself once for
all; in the Mass that same offering is sacramentally made present and the
living Christ is offered unbloodily, without a new immolation. The Church
offers only because she is incorporated into Christ and receives from him what
she offers.

